\documentclass[11pt,a4paper]{ctexart}
\usepackage[margin=2.2cm]{geometry}
\usepackage{xcolor}
\usepackage{enumitem}
\usepackage{amssymb}
\usepackage{titlesec}
\usepackage{hyperref}
\definecolor{navy}{HTML}{1A365D}
\definecolor{ink}{HTML}{1F2933}
\definecolor{muted}{HTML}{66788A}
\definecolor{soft}{HTML}{EEF4FA}
\definecolor{gold}{HTML}{8A5A12}
\hypersetup{colorlinks=true,linkcolor=navy,urlcolor=navy}
\setlist{nosep}
\titleformat{\section}{\Large\bfseries\color{navy}}{}{0pt}{}
\titleformat{\subsection}{\large\bfseries\color{navy}}{}{0pt}{}
\pagestyle{plain}

\title{\textbf{关系从句专项}\\[4pt]
\large 人教版（新课纲版）九年级上册 · 配合 Unit 7 Feel the Rhythm 音乐话题}
\author{初中英语资料站}
\date{2026-11-12}
\begin{document}
\maketitle

\section*{使用说明}
\begin{itemize}[leftmargin=1.4em]
\item 本讲义对应人教版（新课纲版）九年级上册语法要求：关系代词引导的限制性关系从句。
\item 共 4 讲语法、6 级 48 题；建议按“识别 → 选词 → 成分/省略 → 非限制性 → 综合”的顺序完成。
\item 在线版：\url{https://1901151458-prog.github.io/english-materials/relative-clauses-g9/}
\end{itemize}

\section{关系从句是什么}
关系从句由\textbf{关系词}引出，用来修饰前面的名词或代词，被修饰的词叫\textbf{先行词}。因起定语作用，也叫\textbf{定语从句}。
\begin{itemize}[leftmargin=1.4em]
\item I love traditional music \underline{which tells a story}.（先行词 music）
\item The man \underline{who never gives up} is my father.（先行词 the man）
\end{itemize}

\section{选关系代词：that / which / who}
\begin{center}
\begin{tabular}{p{2.4cm}p{2.6cm}p{8.4cm}}
\hline
\textbf{关系代词} & \textbf{指代} & \textbf{例句}\\
\hline
that & 物或人 & I prefer music \underline{that} touches me.\\
which & 物 & The song \underline{which} uses the erhu is my favourite.\\
who & 人 & He is the man \underline{who} never gives up.\\
whose & 所属关系 & The girl \underline{whose} hair is long is my sister.\\
\hline
\end{tabular}
\end{center}
\textbf{选择顺序：}① 先行词是人还是物 → ② 关系代词在从句里作什么成分（主语／宾语／定语）。

\section{作主语还是作宾语：省略规则}
\begin{itemize}[leftmargin=1.4em]
\item 关系代词在从句中作\textbf{宾语}时可以省略：\\
The man \underline{(who)} we met is our teacher. / The music \underline{(which)} he downloaded is popular.
\item 关系代词在从句中作\textbf{主语}时不能省略：\\
The man \underline{who} never gives up is my father.（× 去掉 who）
\end{itemize}
\textbf{判断方法：}把关系代词去掉，看从句是否还缺主语——缺主语不能省；只缺宾语可以省。

\section{非限制性关系从句}
用\textbf{逗号}与主句隔开的补充说明叫非限制性关系从句，删掉后主句仍成立。
\begin{center}
\begin{tabular}{p{2.6cm}p{5.4cm}p{5.4cm}}
\hline
 & \textbf{限制性} & \textbf{非限制性}\\
\hline
逗号 & 无 & 必须有\\
关系词 & that / which / who & which / who（不用 that）\\
省略 & 作宾语可省 & 不可省\\
先行词 & 名词 / 代词 & 名词 / 代词 / 整个主句\\
\hline
\end{tabular}
\end{center}
\begin{itemize}[leftmargin=1.4em]
\item My brother, \underline{who lives in Beijing}, is a doctor.（补充说明先行词）
\item He passed the exam, \underline{which made his parents happy}.（which 指代整个主句）
\end{itemize}


\newpage
\section{分级练习（48 题）}
\noindent 按难度分级排列，先做基础，再做成分与扩展，最后做综合题。

\subsection*{第 1 级 · 基础 · 选关系代词}
\noindent\textcolor{gray}{that / which / who / whose 各自指代什么。}\par\smallskip
\begin{enumerate}[leftmargin=*,itemsep=4pt]
\item I like music \_\_\_ touches my heart.
\begin{enumerate}[label=\Alph*.,leftmargin=2em,itemsep=1pt,topsep=2pt]
\item who
\item which
\item whose
\item where
\end{enumerate}
\item The man \_\_\_ is standing there is my teacher.
\begin{enumerate}[label=\Alph*.,leftmargin=2em,itemsep=1pt,topsep=2pt]
\item which
\item who
\item whose
\item what
\end{enumerate}
\item This is the book \_\_\_ I bought yesterday.
\begin{enumerate}[label=\Alph*.,leftmargin=2em,itemsep=1pt,topsep=2pt]
\item who
\item which
\item where
\item when
\end{enumerate}
\item The girl \_\_\_ hair is long is my sister.
\begin{enumerate}[label=\Alph*.,leftmargin=2em,itemsep=1pt,topsep=2pt]
\item who
\item which
\item whose
\item what
\end{enumerate}
\item The song \_\_\_ uses the erhu is my favourite.
\begin{enumerate}[label=\Alph*.,leftmargin=2em,itemsep=1pt,topsep=2pt]
\item who
\item which
\item whose
\item where
\end{enumerate}
\item I have a friend \_\_\_ can speak French.
\begin{enumerate}[label=\Alph*.,leftmargin=2em,itemsep=1pt,topsep=2pt]
\item which
\item who
\item whose
\item what
\end{enumerate}
\item The house \_\_\_ we live in is very old.
\begin{enumerate}[label=\Alph*.,leftmargin=2em,itemsep=1pt,topsep=2pt]
\item who
\item which
\item whose
\item what
\end{enumerate}
\item The dog \_\_\_ is running is mine.
\begin{enumerate}[label=\Alph*.,leftmargin=2em,itemsep=1pt,topsep=2pt]
\item who
\item which
\item whose
\item where
\end{enumerate}
\end{enumerate}

\subsection*{第 2 级 · 辨析 · 指人还是指物}
\noindent\textcolor{gray}{先看先行词是人还是物，再决定用哪个关系代词。}\par\smallskip
\begin{enumerate}[leftmargin=*,itemsep=4pt]
\item The doctor \_\_\_ treated me is very kind.
\begin{enumerate}[label=\Alph*.,leftmargin=2em,itemsep=1pt,topsep=2pt]
\item which
\item who
\item whose
\item what
\end{enumerate}
\item The chair \_\_\_ I sit on is broken.
\begin{enumerate}[label=\Alph*.,leftmargin=2em,itemsep=1pt,topsep=2pt]
\item who
\item which
\item whose
\item whom
\end{enumerate}
\item The students \_\_\_ passed the exam got prizes.
\begin{enumerate}[label=\Alph*.,leftmargin=2em,itemsep=1pt,topsep=2pt]
\item which
\item who
\item whose
\item what
\end{enumerate}
\item The film \_\_\_ we watched last night was boring.
\begin{enumerate}[label=\Alph*.,leftmargin=2em,itemsep=1pt,topsep=2pt]
\item who
\item which
\item whose
\item what
\end{enumerate}
\item The people \_\_\_ live next door are noisy.
\begin{enumerate}[label=\Alph*.,leftmargin=2em,itemsep=1pt,topsep=2pt]
\item which
\item who
\item whose
\item what
\end{enumerate}
\item The computer \_\_\_ costs too much isn’t good.
\begin{enumerate}[label=\Alph*.,leftmargin=2em,itemsep=1pt,topsep=2pt]
\item who
\item which
\item whose
\item what
\end{enumerate}
\item The teacher \_\_\_ teaches us English is from Canada.
\begin{enumerate}[label=\Alph*.,leftmargin=2em,itemsep=1pt,topsep=2pt]
\item which
\item who
\item whose
\item what
\end{enumerate}
\item The pen \_\_\_ he lost is on the desk.
\begin{enumerate}[label=\Alph*.,leftmargin=2em,itemsep=1pt,topsep=2pt]
\item who
\item which
\item whose
\item what
\end{enumerate}
\end{enumerate}

\subsection*{第 3 级 · 成分 · 作主语还是作宾语}
\noindent\textcolor{gray}{关系代词在从句里作主语还是宾语，直接决定能否省略。}\par\smallskip
\begin{enumerate}[leftmargin=*,itemsep=4pt]
\item The boy \_\_\_ won the game is my brother.
\begin{enumerate}[label=\Alph*.,leftmargin=2em,itemsep=1pt,topsep=2pt]
\item who
\item which
\item whose
\item what
\end{enumerate}
\item The film \_\_\_ we saw is exciting.
\begin{enumerate}[label=\Alph*.,leftmargin=2em,itemsep=1pt,topsep=2pt]
\item who
\item which
\item whose
\item what
\end{enumerate}
\item The people \_\_\_ helped us were very kind.
\begin{enumerate}[label=\Alph*.,leftmargin=2em,itemsep=1pt,topsep=2pt]
\item who
\item which
\item whose
\item what
\end{enumerate}
\item The music \_\_\_ he downloaded is popular.
\begin{enumerate}[label=\Alph*.,leftmargin=2em,itemsep=1pt,topsep=2pt]
\item who
\item which
\item whose
\item what
\end{enumerate}
\item The woman \_\_\_ works in the hospital is a doctor.
\begin{enumerate}[label=\Alph*.,leftmargin=2em,itemsep=1pt,topsep=2pt]
\item who
\item which
\item whose
\item what
\end{enumerate}
\item The story \_\_\_ she told was moving.
\begin{enumerate}[label=\Alph*.,leftmargin=2em,itemsep=1pt,topsep=2pt]
\item who
\item which
\item whose
\item what
\end{enumerate}
\item The students \_\_\_ study hard will succeed.
\begin{enumerate}[label=\Alph*.,leftmargin=2em,itemsep=1pt,topsep=2pt]
\item who
\item which
\item whose
\item what
\end{enumerate}
\item The book \_\_\_ I borrowed is interesting.
\begin{enumerate}[label=\Alph*.,leftmargin=2em,itemsep=1pt,topsep=2pt]
\item who
\item which
\item whose
\item what
\end{enumerate}
\end{enumerate}

\subsection*{第 4 级 · 非限制性 · 逗号 + which}
\noindent\textcolor{gray}{有逗号的补充说明从句：不用 that、不能省、可指代整句。}\par\smallskip
\begin{enumerate}[leftmargin=*,itemsep=4pt]
\item My father, \_\_\_ works in a hospital, is a doctor.
\begin{enumerate}[label=\Alph*.,leftmargin=2em,itemsep=1pt,topsep=2pt]
\item that
\item who
\item which
\item what
\end{enumerate}
\item The book, \_\_\_ I bought last week, is very interesting.
\begin{enumerate}[label=\Alph*.,leftmargin=2em,itemsep=1pt,topsep=2pt]
\item that
\item which
\item who
\item what
\end{enumerate}
\item He passed the exam, \_\_\_ made his parents happy.
\begin{enumerate}[label=\Alph*.,leftmargin=2em,itemsep=1pt,topsep=2pt]
\item that
\item which
\item what
\item who
\end{enumerate}
\item She was late, \_\_\_ was unusual.
\begin{enumerate}[label=\Alph*.,leftmargin=2em,itemsep=1pt,topsep=2pt]
\item that
\item which
\item what
\item who
\end{enumerate}
\item Beijing, \_\_\_ is the capital of China, has a long history.
\begin{enumerate}[label=\Alph*.,leftmargin=2em,itemsep=1pt,topsep=2pt]
\item that
\item which
\item who
\item what
\end{enumerate}
\item Tom, \_\_\_ is my best friend, lives next door.
\begin{enumerate}[label=\Alph*.,leftmargin=2em,itemsep=1pt,topsep=2pt]
\item that
\item which
\item who
\item what
\end{enumerate}
\item The song, \_\_\_ uses the violin, is my favourite.
\begin{enumerate}[label=\Alph*.,leftmargin=2em,itemsep=1pt,topsep=2pt]
\item that
\item which
\item who
\item what
\end{enumerate}
\item He didn’t come, \_\_\_ made us worried.
\begin{enumerate}[label=\Alph*.,leftmargin=2em,itemsep=1pt,topsep=2pt]
\item that
\item which
\item what
\item who
\end{enumerate}
\end{enumerate}

\subsection*{第 5 级 · 省略 · 什么时候能省}
\noindent\textcolor{gray}{作宾语可省，作主语不可省 —— 判断题专练。}\par\smallskip
\begin{enumerate}[leftmargin=*,itemsep=4pt]
\item The boy \_\_\_ broke the window is Tom. 此空能否省略？
\begin{enumerate}[label=\Alph*.,leftmargin=2em,itemsep=1pt,topsep=2pt]
\item 可以省略
\item 不能省略
\item 只能换成 that
\item 只能换成 which
\end{enumerate}
\item The pen \_\_\_ he used is mine. 此空能否省略？
\begin{enumerate}[label=\Alph*.,leftmargin=2em,itemsep=1pt,topsep=2pt]
\item 可以省略
\item 不能省略
\item 只能换成 who
\item 必须保留 that
\end{enumerate}
\item The man \_\_\_ we met is my uncle. 此空能否省略？
\begin{enumerate}[label=\Alph*.,leftmargin=2em,itemsep=1pt,topsep=2pt]
\item 可以省略
\item 不能省略
\item 只能换成 which
\item 必须用 whose
\end{enumerate}
\item The woman \_\_\_ works here is my aunt. 此空能否省略？
\begin{enumerate}[label=\Alph*.,leftmargin=2em,itemsep=1pt,topsep=2pt]
\item 可以省略
\item 不能省略
\item 只能换成 which
\item 必须用 whose
\end{enumerate}
\item The music \_\_\_ I like is soft. 此空能否省略？
\begin{enumerate}[label=\Alph*.,leftmargin=2em,itemsep=1pt,topsep=2pt]
\item 可以省略
\item 不能省略
\item 只能换成 who
\item 必须保留 which
\end{enumerate}
\item The students \_\_\_ study hard will succeed. 此空能否省略？
\begin{enumerate}[label=\Alph*.,leftmargin=2em,itemsep=1pt,topsep=2pt]
\item 可以省略
\item 不能省略
\item 只能换成 which
\item 必须用 whose
\end{enumerate}
\item The book \_\_\_ she bought is new. 此空能否省略？
\begin{enumerate}[label=\Alph*.,leftmargin=2em,itemsep=1pt,topsep=2pt]
\item 可以省略
\item 不能省略
\item 只能换成 who
\item 必须保留 that
\end{enumerate}
\item The girl \_\_\_ won the prize is my cousin. 此空能否省略？
\begin{enumerate}[label=\Alph*.,leftmargin=2em,itemsep=1pt,topsep=2pt]
\item 可以省略
\item 不能省略
\item 只能换成 which
\item 必须用 whose
\end{enumerate}
\end{enumerate}

\subsection*{第 6 级 · 综合 · 语境提升}
\noindent\textcolor{gray}{最高级、only、everything、非限制性等常考陷阱混合辨析。}\par\smallskip
\begin{enumerate}[leftmargin=*,itemsep=4pt]
\item This is the best film \_\_\_ I have ever seen.
\begin{enumerate}[label=\Alph*.,leftmargin=2em,itemsep=1pt,topsep=2pt]
\item who
\item that
\item what
\item whose
\end{enumerate}
\item The man, \_\_\_ I met yesterday, is a writer.
\begin{enumerate}[label=\Alph*.,leftmargin=2em,itemsep=1pt,topsep=2pt]
\item that
\item who
\item which
\item what
\end{enumerate}
\item I don’t like music \_\_\_ is too loud.
\begin{enumerate}[label=\Alph*.,leftmargin=2em,itemsep=1pt,topsep=2pt]
\item who
\item which
\item whose
\item where
\end{enumerate}
\item He is the only person \_\_\_ I can trust.
\begin{enumerate}[label=\Alph*.,leftmargin=2em,itemsep=1pt,topsep=2pt]
\item who
\item that
\item which
\item what
\end{enumerate}
\item The village \_\_\_ we visited last year has changed a lot.
\begin{enumerate}[label=\Alph*.,leftmargin=2em,itemsep=1pt,topsep=2pt]
\item who
\item which
\item whose
\item where
\end{enumerate}
\item She told us everything \_\_\_ she knew.
\begin{enumerate}[label=\Alph*.,leftmargin=2em,itemsep=1pt,topsep=2pt]
\item who
\item that
\item which
\item what
\end{enumerate}
\item The house \_\_\_ we live in is very old.
\begin{enumerate}[label=\Alph*.,leftmargin=2em,itemsep=1pt,topsep=2pt]
\item who
\item which
\item whose
\item what
\end{enumerate}
\item He passed the exam, \_\_\_ surprised us all.
\begin{enumerate}[label=\Alph*.,leftmargin=2em,itemsep=1pt,topsep=2pt]
\item that
\item which
\item what
\item who
\end{enumerate}
\end{enumerate}

\newpage
\section*{答案与解析}
\small
\begin{enumerate}[leftmargin=1.6em,itemsep=2pt]
\item \textbf{R1-1 B}：先行词 music 是物，关系代词在从句中作主语，用 which（也可用 that）。
\item \textbf{R1-2 B}：先行词 the man 是人，关系代词作主语，用 who。
\item \textbf{R1-3 B}：先行词 the book 是物，用 which（作宾语，也可用 that 或省略）。
\item \textbf{R1-4 C}：表示「她的头发」，关系代词作定语，用所有格 whose。
\item \textbf{R1-5 B}：先行词 the song 是物，作主语，用 which（教材例句）。
\item \textbf{R1-6 B}：先行词 a friend 是人，作主语，用 who。
\item \textbf{R1-7 B}：先行词 the house 是物，作介词 in 的宾语，用 which。
\item \textbf{R1-8 B}：先行词 the dog 是动物（物），用 which。
\item \textbf{R2-1 B}：先行词 the doctor 是人，作主语，用 who。
\item \textbf{R2-2 B}：先行词 the chair 是物，用 which（也可用 that）。
\item \textbf{R2-3 B}：先行词 the students 是人，作主语，用 who。
\item \textbf{R2-4 B}：先行词 the film 是物，作宾语，用 which。
\item \textbf{R2-5 B}：先行词 the people 是人，作主语，用 who。
\item \textbf{R2-6 B}：先行词 the computer 是物，作主语，用 which。
\item \textbf{R2-7 B}：先行词 the teacher 是人，作主语，用 who。
\item \textbf{R2-8 B}：先行词 the pen 是物，作宾语，用 which。
\item \textbf{R3-1 A}：先行词 the boy 是人；关系代词在从句中作主语（won 的主语），用 who。
\item \textbf{R3-2 B}：先行词 the film 是物；which 在从句中作 saw 的宾语。
\item \textbf{R3-3 A}：先行词 the people 是人；关系代词作 helped 的主语，用 who。
\item \textbf{R3-4 B}：先行词 the music 是物；which 在从句中作 downloaded 的宾语。
\item \textbf{R3-5 A}：先行词 the woman 是人；关系代词作 works 的主语，用 who。
\item \textbf{R3-6 B}：先行词 the story 是物；which 作 told 的宾语。
\item \textbf{R3-7 A}：先行词 the students 是人；关系代词作 study 的主语，用 who。
\item \textbf{R3-8 B}：先行词 the book 是物；which 作 borrowed 的宾语。
\item \textbf{R4-1 B}：逗号隔开的非限制性从句，指人用 who，不能用 that。
\item \textbf{R4-2 B}：非限制性从句指物用 which，不能用 that。
\item \textbf{R4-3 B}：which 指代前面整个主句（他通过考试这件事），不能用 that 或 what。
\item \textbf{R4-4 B}：which 指代「她迟到」这件事，非限制性从句指代整句只能用 which。
\item \textbf{R4-5 B}：非限制性从句指物（Beijing 作地点名词），用 which，不用 that。
\item \textbf{R4-6 C}：非限制性从句指人用 who，不能用 that。
\item \textbf{R4-7 B}：非限制性从句指物用 which，不能用 that。
\item \textbf{R4-8 B}：which 指代「他没来」这件事，非限制性从句指代整句只用 which。
\item \textbf{R5-1 B}：关系代词在从句中作主语（broke 的主语），作主语不能省略。
\item \textbf{R5-2 A}：关系代词作 used 的宾语，作宾语时可以省略。
\item \textbf{R5-3 A}：关系代词作 met 的宾语，作宾语可省略。
\item \textbf{R5-4 B}：关系代词作 works 的主语，作主语不能省略。
\item \textbf{R5-5 A}：关系代词作 like 的宾语，作宾语可省略。
\item \textbf{R5-6 B}：关系代词作 study 的主语，作主语不能省略。
\item \textbf{R5-7 A}：关系代词作 bought 的宾语，作宾语可省略。
\item \textbf{R5-8 B}：关系代词作 won 的主语，作主语不能省略。
\item \textbf{R6-1 B}：先行词被最高级 the best 修饰时，关系代词常用 that（也可省略，因作宾语）。
\item \textbf{R6-2 B}：逗号隔开的非限制性从句指人，用 who，不能用 that。
\item \textbf{R6-3 B}：先行词 music 是物，关系代词作主语，用 which。
\item \textbf{R6-4 B}：先行词被 the only 修饰时，关系代词常用 that。
\item \textbf{R6-5 B}：从句缺 visited 的宾语，用关系代词 which；若从句不缺宾语才用 where。
\item \textbf{R6-6 B}：先行词是 everything 等不定代词时，关系代词通常用 that。
\item \textbf{R6-7 B}：关系代词作介词 in 的宾语，指物用 which。
\item \textbf{R6-8 B}：非限制性从句指代整个主句，只能用 which。
\end{enumerate}

\end{document}
